\emph{The bounce lemma (recalled; filed alongside, synthesis.bkl-bounce-invariant-holonomy-and-bound).} With $d = 1 + 2p_{a}$ (positive, since $p_{a} \geq -\tfrac13$ on the constraint variety), $\sum_{i} p'_{i} = (-p_{a} + \sum_{b \neq a} p_{b} + 4p_{a})/d = (1 + 2p_{a})/d = 1$, and $\sum_{i} p_{i}'^{2} = \bigl(p_{a}^{2} + \sum_{b\neq a}(p_{b} + 2p_{a})^{2}\bigr)/d^{2} = \bigl(1 + 4p_{a}(1 - p_{a}) + 8p_{a}^{2}\bigr)/d^{2} = (1 + 2p_{a})^{2}/d^{2} = 1$, using $\sum_{b \ne a} p_{b} = 1 - p_{a}$ and $\sum_{b\neq a}p_{b}^{2} = 1 - p_{a}^{2}$. On $p(u)$, $a$ is the first axis and $d = (u^{2} - u + 1)/(u^{2} + u + 1)$; with $D' = u^{2} - u + 1 = 1 + (u - 1) + (u - 1)^{2}$ one gets $p'_{a} = u/D' = p_{2}(u - 1)$, $p'_{b} = (1 - u)/D' = p_{1}(u - 1)$ and $p'_{c} = u(u - 1)/D' = p_{3}(u - 1)$. For $0 < w < 1$ the triple $p(w)$ is a permutation of $p(1/w)$ (multiply numerators and denominator by $w^{-2}$), so after a step to $u - 1 < 1$ the ordered triple is $p(1/(u - 1))$.

\emph{(i).} Assume the approach follows the transition law. The coefficient of $B_{a}$ goes from $p_{1}(u) = -u/(u^{2} + u + 1)$ to $p_{2}(u - 1) = u/(u^{2} - u + 1)$, a change $j(u) = u/(u^{2} + u + 1) + u/(u^{2} - u + 1)$. Both summands have derivative numerator $1 - u^{2} \leq 0$ for $u \geq 1$, so on $[1, 2]$, $j(u) \geq j(2) = \tfrac27 + \tfrac23 = \tfrac{20}{21}$. An era starting at $u = k + x$ ($k \geq 1$, $0 \leq x < 1$) runs through $k + x, k - 1 + x, \ldots, 1 + x$, and its last epoch has $1 \leq u < 2$. For irrational initial $u$ no epoch value is an integer, every era ends with $u - 1 \in (0, 1)$ and the next starts at $1/(u - 1) > 1$; this is the continued-fraction expansion, which is infinite, so there are infinitely many eras and infinitely many jumps of size at least $20/21$ in a fixed-frame coefficient. A Cauchy sequence has eventually small jumps, so the sequence $\widehat\Omega_{0}(p^{(n)})$ is not Cauchy in any norm on the coefficients.

\emph{(ii).} For rational $u$ the steps $u \mapsto u - 1$ and $u \mapsto 1/(u - 1)$ are Euclid's algorithm on numerator and denominator: within an era the integer part decreases by one per bounce, and an era change replaces the fractional part $x$ by $1/x$, which has a smaller denominator. So the orbit reaches an integer and then $u = 1$, after $a_{0} + \cdots + a_{k} - 1$ bounces, and $p(1) = (-\tfrac13, \tfrac23, \tfrac23)$, the Schwarzschild exponents of Lemma \claimref{interior.kasner-form-interior}. There $d = \tfrac13$, $p'_{a} = 1$ and $p'_{b} = p'_{c} = (\tfrac23 - \tfrac23)/\tfrac13 = 0$. The triple $(1, 0, 0)$ has no negative exponent, so the law does not apply and the orbit stops; it has two zero exponents, so it violates $p_{i} \neq 0$, and the two terminal triples are exactly the repeated-exponent triples of Proposition \claimref{sheaf.collision-census}. The limiting values at $(1, 0, 0)$ are multiples of the one boost $B_{a}$, which spans a one-dimensional abelian algebra. At every other point of the constraint variety all three exponents are nonzero: by Proposition \claimref{sheaf.collision-census}(ii), if one exponent vanishes the triple is $(0, 0, 1)$ up to permutation, so two nonzero exponents force the third. The limiting values then include nonzero multiples of all three boosts $B_{1}, B_{2}, B_{3}$, and these generate exactly $\Lambda^{2}W \cong \mathfrak{so}(1,3)$ by Proposition \claimref{interior.boosts-generate-exactly-levi-factor}. (Two boosts alone generate only a three-dimensional algebra, a copy of $\mathfrak{so}(1,2)$; the third boost is what gives six dimensions.) The check computes these closures in $\Cl(1,3)$ and verifies the count of bounces against the continued fraction for eight rational values.

\emph{(iii).} The absence of nonconstant continuous invariants is classical for this map (the Gauss-map dynamics of the BKL parameter); the argument is given for completeness. The unordered triple is the root set of $t^{3} - t^{2} - e_{3}$ (as $e_{1} = 1$ and $e_{2} = \tfrac12(e_{1}^{2} - \sum p_{i}^{2}) = 0$), so a function of the unordered triple is a function of $e_{3}$, and on $u \geq 1$, $e_{3}(u) = -u^{2}(1 + u)^{2}/(1 + u + u^{2})^{3}$, whose derivative has numerator $u(1 + u)(1 + 2u)(u + 2)(u - 1)(1 + u + u^{2})^{2} > 0$ for $u > 1$: $e_{3}$ is a continuous bijection from $[1, \infty)$ onto $[-\tfrac{4}{27}, 0)$. Let $f$ be continuous and invariant. By (ii), every rational $r \geq 1$ is carried by finitely many bounces to $u = 1$, so $f(e_{3}(r)) = f(-\tfrac{4}{27})$; rationals are dense in $[1, \infty)$ and $f\circ e_{3}$ is continuous, so $f\circ e_{3}$ is constant on $[1, \infty)$, and $f$ is constant (at $e_{3} = 0$ by continuity). A bounce from $u$ preserves $e_{3}$ exactly when $e_{3}(u) = e_{3}(u')$ with $u'$ the next ordered parameter: for $u \geq 2$, $u' = u - 1 \neq u$; for $1 < u < 2$, $u' = 1/(u - 1)$, equal to $u$ exactly when $u^{2} - u - 1 = 0$; injectivity of $e_{3}$ on $[1,\infty)$ does the rest, and from $u = 1$ the bounce leaves the domain ($e_{3}$ goes from $-\tfrac{4}{27}$ to $0$). The wall counterterm's coefficient is $8p_{1}p_{2}p_{3} = 8e_{3}$ (Theorem \claimref{dark-energy.wall-counterterm}(ii)), so it changes at every bounce except from the golden epoch. The limiting connection being the one of Theorem \claimref{interior.weighted-limit-theorem}(3) on each epoch is the hypothesis that each epoch is admissible; the theorem itself is not extended here.
